
\section{泛函方法}
\subsection{一般的路径积分方法}
从量子理论的正则对易关系出发，其有正则``坐标'' $Q_a$ 以及共轭``动量'' $P_b$，
满足正则对易关系\footnote{%
  我们在这里隐含假定了任何第一类约束通过选择规范被消除了，
  而剩下的任何第二类约束通过将约束自由度
  写成非约束 $Q_a$ 和 $P_a$ 的形式被``解出''了，
  就像 \S5.2.3 中做的那样。
  路径积分方法对约束系统的直接应用被 Faddeev 所描述
  （参见讨论~2）。}：
\begin{subequations}\label{eq:CCR-gen}
\begin{align}
  [Q_a,\,P_b] &= i\,\delta_{ab}\,,\label{eq:QP}\\
  [Q_a,\,Q_b] &= [P_a,\,P_b] = 0\,.\label{eq:QQ-PP}
\end{align}
\end{subequations}
指标 $a$ 由空间位置 $\mathbf{x}$
和一个离散的 Lorentz 及种类指标 $n$ 构成：
\begin{equation}\label{eq:field-index}
  Q_a \;\longleftrightarrow\; Q_n(\mathbf{x})\,,\qquad
  P_a \;\longleftrightarrow\; P_n(\mathbf{x})\,,\qquad
  \delta_{ab} \;\longleftrightarrow\;
    \delta_{nm}\,\delta^{(3)}(\mathbf{x}-\mathbf{y})\,.
\end{equation}
这些算符是``Schr\"odinger-绘景''算符，取在固定时刻（例如 $t=0$）。
\medskip


既然 $Q_a$ 和$P_b$全部对易，可以构造本征值为 $q_a$和$p_a$ 的同时本征态：
\begin{subequations}
  \begin{align}
    &Q_a\ket{q} = q_a\ket{q}\\
    &P_a\ket{p} = p_a\ket{p}
  \end{align}
\end{subequations}
这些本征矢是完备且正交的
\begin{subequations}
  \begin{align}
    &\braket{q'|q} = \prod_a \delta(q'_a - q_a)
    \equiv \delta(q'-q)\\
    &\mathbf{1} = \int\prod_a dq_a\;\ket{q}\bra{q}\\
    &\braket{p'|p} = \delta(p'-p)\\
    &\mathbf{1} = \int\prod_a dp_a\;\ket{p}\bra{p}
  \end{align}
\end{subequations}
从不对易的关系 (\ref{eq:QP})，可以得出这两组完备本征态之间的标量积
\begin{equation}\label{eq:qp-inner}
  \braket{q|p} = \prod_a \frac{1}{\sqrt{2\pi}}\,
    \exp(i\,q_a\,p_a)\,.
\end{equation}
算符$Q_a$ 和$P_b$全部的时间依赖性都由时间演化算符确定，由$H$总的哈密顿量确定；它们
严格定义在 Heisenberg 绘景中
\begin{subequations}\label{eq:Heis-ops}
  \begin{align}
    &Q_a(t) \equiv e^{iHt}\,Q_a\,e^{-iHt}\\
    &P_a(t) \equiv e^{iHt}\,P_a\,e^{-iHt}
  \end{align}
\end{subequations}
以及本征态的时间依赖性：
\begin{subequations}\label{eq:Heis-states}
  \begin{align}
  &\ket{q;t} = e^{iHt}\ket{q}\\
  &\ket{p;t} = e^{iHt}\ket{p}
  \end{align}
\end{subequations}
可以证明 $\ket{q;t}$ 是 $Q_a(t)$ 本征值为 $q_a$ 的本征态($\ket{p}$同理):
\[Q_a(t)\ket{q;t}=e^{iHt}\,Q_a\,e^{-iHt}e^{iHt}\ket{q}=q_ae^{iHt}\ket{q}=q_a\ket{q;t}\]
这些态满足
\begin{subequations}
  \begin{align}
  &\braket{q';t|q;t}=\delta(q'-q)\,,\qquad
  \mathbf{1}=\int\prod_a dq_a\;\ket{q;t}\bra{q;t}\\
  &\braket{p';t|p;t}=\delta(p'-p)\,,\qquad
  \mathbf{1}=\int\prod_a dp_a\;\ket{p;t}\bra{p;t}
  \end{align}
\end{subequations}
以及
\begin{equation}\label{eq:Heis-qp}
  \braket{q;t|p;t}
  =\prod_a\frac{1}{\sqrt{2\pi}}\,\exp(i\,q_a\,p_a)\,.
\end{equation}
我们核心的动力学问题是计算标量积
$\braket{q';t'|q;t}$（$t'>t$,
\begin{equation}\label{eq:Feynman_k}
  \braket{q';t'|q;t}=\braket{q'|e^{i(t'-t)H(Q(t),P(t))}|q}
\end{equation}
然而直接计算是十分不易的，原因如下：
\begin{enumerate}
  \item $(t'-t)$是一个有限间隔，也就是说$H$并不是线性的；
  \item $H=H(Q(t),P(t))$是$Q(t)$和$P(t)$的泛函，是一个算符，并不能直接计算。
\end{enumerate}
上述问题解决的办法：
\begin{enumerate}
  \item 将时间切片，对于无穷小的间隔，$H$是线性的；
  \item 在其中一段时间间隔内插入完备关系$\mathbf{1} = \int\prod_a dp_a\;\ket{p}\bra{p}$，并规定$H$中的$Q(t)$在$P(t)$的左边这样的标准形式，对于不是标准形式的，可以利用对易关系(\refeq{eq:QP})替换为标准形式，比如给定哈密顿量中一项的形式为 $P_a\,Q_b\,P_c$，将其重写为
        \[
          P_a\,Q_b\,P_c = Q_b\,P_a\,P_c - i\,\delta_{ab}\,P_c\,.
        \]
\end{enumerate}
这样的处理方式，使得我们可以计算标量积(\refeq{eq:Feynman_k}).
\medskip


取一个小间隔，当 $t'$ 和 $t$ 无限接近时，
即 $t'=\tau+d\tau$、$t=\tau$，利用 (\ref{eq:Heis-states})：
\begin{equation}\label{eq:inf-step-0}
  \braket{q';\tau{+}d\tau\,|\,q;\tau}
  =\bra{q';\tau}\exp(-iH\,d\tau)\ket{q;\tau}\,.
\end{equation}
在 $P$-本征态下展开 $\ket{p;\tau}$：
\begin{align}
  \braket{q';\tau{+}d\tau\,|\,q;\tau}
  &= \int\prod_a dp_a\;
    \bra{q';\tau}\exp\bigl(-iH(Q(\tau),P(\tau))\,d\tau\bigr)
    \ket{p;\tau}\braket{p;\tau|q;\tau}
    \notag\\[6pt]
  &= \int\prod_a\frac{dp_a}{2\pi}\;
    \exp\!\biggl\{
      -iH(q',p)\,d\tau
      +i\sum_a(q'_a-q_a)\,p_a
    \biggr\}\,,
    \label{eq:single-step}
\end{align}
其中每个 $p_a$ 从 $-\infty$ 积到 $+\infty$。
\medskip


回到有限间隔$[t,t']$中，为了计算 $t<t'$ 的 $\braket{q';t'|q;t}$，
将从 $t$ 到 $t'$ 的时间间隔分解成
$t,\,\tau_1,\,\tau_2,\,\cdots,\,\tau_N,\,t'$，其中
\begin{equation}\label{eq:dt}
  \tau_{k+1}-\tau_k = d\tau = \frac{t'-t}{N+1}\,,
\end{equation}
并对每个时刻 $\tau_k$ 的态 $\ket{q;\tau_k}$
的完全集求和：
\begin{equation}\label{eq:slicing}
  \braket{q';t'|q;t}
  =\int dq_1\cdots dq_N\;
    \braket{q';t'|q_N;\tau_N}
    \braket{q_N;\tau_N|q_{N-1};\tau_{N-1}}
    \cdots
    \braket{q_1;\tau_1|q;t}\,.
\end{equation}
利用 (\ref{eq:single-step})，
令 $q_0\equiv q$、$q_{N+1}\equiv q'$，得到
\begin{equation}\label{eq:PI-discrete}
  \begin{aligned}
      \braket{q';t'|q;t}
  =\int\biggl[\prod_{k=1}^{N}\prod_a dq_{k,a}\biggr]
    &\biggl[\prod_{k=0}^{N}\prod_a
      \frac{dp_{k,a}}{2\pi}\biggr]\\\times
  &\exp\!\biggl[
    i\sum_{k=1}^{N}\biggl\{
      \sum_a(q_{k+1,a}-q_{k,a})\,p_{k,a}
      -H(q_k+1,p_{k})\,d\tau
    \biggr\}
  \biggr]\,.
  \end{aligned}
\end{equation}

\medskip\noindent
我们将$q$和$p$的指标$k$的依赖性放到时间中，定义光滑插值函数 $q_a(\tau)$ 和 $p_a(\tau)$，
使得 $q_a(\tau_k)\equiv q_{k,a}$、$p_a(\tau_k)\equiv p_{k,a}$。
在极限 $d\tau\to 0$（即 $N\to\infty$）下，
(\ref{eq:PI-discrete}) 中指数的求和变成对 $\tau$ 的积分：
\begin{align}
  &\sum_{k=1}^{N}\biggl\{
    \sum_a(q_{k+1,a}-q_{k,a})\,p_{k,a}
    -H(q_k+1,p_{k})\,d\tau
  \biggr\}
  \notag\\[4pt]
  &\quad=\sum_{k=1}^{N}\biggl\{
    \sum_a\dot{q}_a(\tau_k)\,p_a(\tau_k)
    -H\bigl(q(\tau_k),p(\tau_k)\bigr)
  \biggr\}d\tau + \mathscr{O}(d\tau^2)
  \notag\\[4pt]
  &\quad\longrightarrow\;
  \int_{t}^{t'}\!d\tau\;\biggl\{
    \sum_a\dot{q}_a(\tau)\,p_a(\tau)
    -H\bigl(q(\tau),p(\tau)\bigr)
  \biggr\}\,.
  \label{eq:cont-limit}
\end{align}
在这操作下，指标$k$全部换为时间的依赖性，函数 $q_a(\tau)$、$p_a(\tau)$ 的路径积分测度为
\begin{equation}\label{eq:measure}
  \int \mathcal{D}q\frac{\mathcal{D}p}{2\pi}\cdots=\int\prod_{\tau,a}dq_a(\tau)
    \prod_{\tau,b}\;\frac{dp_b(\tau)}{2\pi}\;\cdots
  \;\equiv\;
  \lim_{d\tau\to 0}\int
    \prod_{k,a}dq_{k,a}
    \prod_{k,b}\frac{dp_{k,b}}{2\pi}\;\cdots
\end{equation}
那么 (\ref{eq:PI-discrete}) 变成\textbf{相空间路径积分}：
\begin{equation}\label{eq:PI-phase}
  \begin{aligned}
  \braket{q';t'|q;t}
  =&\int_{\substack{q_a(t)=q_a\\q_a(t')=q'_a}}
    \prod_{\tau,a}dq_a(\tau)
    \prod_{\tau,b}\frac{dp_b(\tau)}{2\pi}\\
    &\times\exp\!\biggl[
    i\int_{t}^{t'}\!d\tau\;\biggl\{
      \sum_a\dot{q}_a(\tau)\,p_a(\tau)
      -H\bigl(q(\tau),p(\tau)\bigr)
    \biggr\}
  \biggr]\\
  =&\int_{\substack{q_a(t)=q_a\\q_a(t')=q'_a}}
    \mathcal{D}q\frac{\mathcal{D}p}{2\pi}\exp\!\biggl[
    i\int_{t}^{t'}\!d\tau\;\biggl\{
      \sum_a\dot{q}_a(\tau)\,p_a(\tau)
      -H\bigl(q(\tau),p(\tau)\bigr)
    \biggr\}
  \biggr]\\
  \end{aligned}
\end{equation}
这称为路径积分，是对所有从 $\tau=t$ 时的 $q$
到 $\tau=t'$ 时的 $q'$ 的路径 $q(\tau)$ 积分，
并对所有 $p(\tau)$ 积分。
\medskip



路径积分形式体系不仅允许我们计算跃迁振幅
$\braket{q';t'|q;t}$，
还可以推广至，插入任意算符$\mathcal{O}\bigl(P(\tau),Q(\tau)\bigr)$
的矩阵元

\medskip\noindent
对于无穷小的间隔的矩阵元，插入算符$\mathcal{O}\bigl(P(\tau),Q(\tau)\bigr)$，这里规定\text{\color{blue}所有的 $P$ 处在左边而所有的 $Q$ 处在右边}，那么有：
\begin{align}
  &\braket{q';\tau{+}d\tau\,|\mathcal{O}\bigl(P(\tau),Q(\tau)\bigr)|\,q;\tau}\\
  &= \int\prod_a dp_a\;
    \bra{q';\tau}\exp\bigl(-iH(Q(\tau),P(\tau))\,d\tau\bigr)
    \ket{p;\tau}\braket{p;\tau|\mathcal{O}\bigl(P(\tau),Q(\tau)\bigr)|q;\tau}
    \notag\\[6pt]
  &= \int\prod_a\frac{dp_a}{2\pi}\mathcal{O}\bigl(p(\tau),q(\tau)\bigr)\;
    \exp\!\biggl\{
      -iH(q',p)\,d\tau
      +i\sum_a(q'_a-q_a)\,p_a
    \biggr\}\,,
    \label{eq:single-step}
\end{align}
回到任意有限间隔的矩阵元，\color{blue} 由于小间隔是按时间切片了，所以插入的算符理应有一定的顺序;例如$t_A\in (t_{k+1},t'_k)$，那么在矩阵元$\braket{q_{k+1};t_{k+1}|q_k;t_k}$之间插入的算符之能是$\mathcal{O}_A \bigl(P(t_A),Q(t_A)\bigr)$ \normalcolor;那么对于任意时间顺序的算符，可以利用\text{编时算符}，使其有：
\[T:t'>t_A>t_B>\dots>t\]
这是唯一的，所以路径积分变成了
\begin{equation}\label{eq:PI-field}
\begin{aligned}
  &\bra{q',t'}T\bigl\{
    \mathcal{O}_A[\,p(t_A),q(t_A)\,]\;
    \mathcal{O}_B[\,p(t_B),q(t_B)\,]\;\cdots
  \bigr\}\ket{q,t}\\[6pt]
  &\;=\int_{\substack{q_a(t)=q_a\\q_a(t')=q'_a}}
    \prod_{\tau,a}dq_a(\tau)\;
    \prod_{\tau,b}
      \frac{dp_b(\tau)}{2\pi}
    \mathcal{O}_A\bigl[\,p(t_A),q(t_A)\,\bigr]\;
    \mathcal{O}_B\bigl[\,p(t_B),q(t_B)\,\bigr]\;\cdots\\
  &\;\quad\times
    \exp\!\biggl[
      i\int_{t}^{t'}\!d\tau\;\biggl\{
        \sum_n\dot{q}_a(\tau)\,
          p_b(\tau)
        -H\bigl[\,q(\tau),p(\tau)\,\bigr]
      \biggr\}
    \biggr]\\
    &\;=\int_{\substack{q_a(t)=q_a\\q_a(t')=q'_a}}
    \mathcal{D}q\frac{\mathcal{D}p}{2\pi}
    \mathcal{O}_A\bigl[\,p(t_A),q(t_A)\,\bigr]\;
    \mathcal{O}_B\bigl[\,p(t_B),q(t_B)\,\bigr]\;\cdots\\
  &\;\quad\times
    \exp\!\biggl[
      i\int_{t}^{t'}\!d\tau\;\biggl\{
        \sum_n\dot{q}_a(\tau)\,
          p_b(\tau)
        -H\bigl[\,q(\tau),p(\tau)\,\bigr]
      \biggr\}
    \biggr]\\
    &\;=\int_{\substack{q_a(t)=q_a\\q_a(t')=q'_a}}
    \mathcal{D}q\frac{\mathcal{D}p}{2\pi}
    \mathcal{O}_A\bigl[\,p(t_A),q(t_A)\,\bigr]\;
    \mathcal{O}_B\bigl[\,p(t_B),q(t_B)\,\bigr]\;e^{iI(q,p)}
\end{aligned}
\end{equation}
其中约定作用量泛函是：
\begin{equation}\label{eq:action-func}
  I[q,p] \equiv \int_{t}^{t'}\!d\tau\;\biggl\{
    \sum_a\dot{q}_a(\tau)\,p_a(\tau)
    -H\bigl(q(\tau),p(\tau)\bigr)
  \biggr\}
\end{equation}
% ─────────────────────────────────────────────
% =============================================================
%   讨论 1：运动方程的涌现与接触项
% =============================================================

\topictitle{***}
\text{\color{blue}或许该强调一下}
在推导路径积分的过程中，算符 $Q_a,P_b$
被替换为 c-数积分变量 $q_a(\tau),p_a(\tau)$。
算符形式体系中的两大核心结构似乎因此消失了：
\begin{itemize}
  \item[(A)] Heisenberg 运动方程
    $\dot{Q}_a=\partial H/\partial P_a$，
    $\dot{P}_a=-\partial H/\partial Q_a$；
  \item[(B)] 正则对易关系
    $[Q_a,\,P_b]=i\,\delta_{ab}$。
\end{itemize}
讨论的核心是：\textbf{完全从路径积分内部，
把 (A) 和 (B) 都找回来}，
从而确认两种表述的等价性。


\bigskip
\medskip\noindent
出发点是一个简单的事实：被积函数对自由积分变量的全导数，积分后为零。 具体地说：
\begin{itemize}
  \item \textbf{$p_a(\tau)$}：
    在\textbf{所有}时刻 $\tau\in[t,t']$
    都是完全自由的积分变量
    （积分区间 $(-\infty,+\infty)$，无边界条件）。
    因此对任意 $t_A\in[t,t']$，
    分部积分
    \begin{equation}\label{eq:IBP-p}
      \int_{-\infty}^{+\infty}\!dp_a(t_A)\;
      \frac{\partial}{\partial p_a(t_A)}[\,\cdots\,]=0
    \end{equation}
    总是成立的。

  \item \textbf{$q_a(\tau)$}：
    在端点 $\tau=t$ 和 $\tau=t'$ 处
    被边界条件 $q_a(t)=q_a$、$q_a(t')=q'_a$
    所\textbf{固定}（不是积分变量）。
    只有在\textbf{中间}时刻 $t<\tau<t'$，
    $q_a(\tau)$ 才是自由积分变量。
    因此分部积分
    \begin{equation}\label{eq:IBP-q}
      \int_{-\infty}^{+\infty}\!dq_a(t_A)\;
      \frac{\partial}{\partial q_a(t_A)}[\,\cdots\,]=0
    \end{equation}
    仅在 $t<t_A<t'$（\textbf{严格不等号}）时成立。
\end{itemize}
这是我们唯一的出发点。
但要利用它推出物理结论，
需要知道
$\frac{\delta e^{iI}}{\delta p_a(t_A)}$以及$\frac{\delta e^{iI}}{\delta q_a(t_A)}$ 到底等于什么?依次考察这两个量：
\[\begin{aligned}
  \frac{\delta e^{iI}}{\delta p_a(t_A)}=
  ie^{iI}\frac{\delta I}{\delta p_a(t_A)}&=ie^{iI}\int_{t}^{t'}d\tau \Biggl\{ \frac{\delta \sum_{b}\dot{q}_b(\tau) p_{b}(\tau)}{\delta p_a(t_A)}-\frac{\delta H(q(\tau),p(\tau))}{\delta p_a(t_A)}\Biggr\}
  \\
  \;&=ie^{iI}\int_{t}^{t'}d\tau \Biggl\{ \sum_{b}\frac{\delta p_{b}(\tau) }{\delta p_a(t_A)}\dot{q}_b(\tau)-\frac{\partial H(q(\tau),p(\tau))}{\partial p_b(\tau)}\frac{\delta p_b(\tau)}{\delta p_a(t_A)}\Biggr\}
  \\
  \;&=ie^{iI}\int_{t}^{t'}d\tau \Biggl\{ \sum_{b}\delta_{ab}\delta(\tau-t_A)\dot{q}_b(\tau)-\frac{\partial H(q(\tau),p(\tau))}{\partial p_b(\tau)}\delta_{ab}\delta(\tau-t_A)\Biggr\}
  \\
  \;&=ie^{iI} \Biggl\{ \dot{q}_a(t_A)-\frac{\partial H(q(t_A),p(t_A))}{\partial p_a(t_A)}\Biggr\}
\end{aligned}\]
以及
\[\begin{aligned}
  \frac{\delta e^{iI}}{\delta q_a(t_A)}=
  ie^{iI}\frac{\delta I}{\delta q_a(t_A)}&=ie^{iI}\int_{t}^{t'}d\tau \Biggl\{ \frac{\delta \sum_{b}\dot{q}_b(\tau) p_{b}(\tau)}{\delta q_a(t_A)}-\frac{\delta H(q(\tau),p(\tau))}{\delta q_a(t_A)}\Biggr\}
  \\
  \;&=ie^{iI}\int_{t}^{t'}d\tau \Biggl\{ \frac{\delta \sum_{b}\frac{d(q_b(\tau)p_b(\tau))}{dt}-q_b(\tau)\dot{p}_b(\tau)}{\delta q_a(t_A)}-\frac{\delta H(q(\tau),p(\tau))}{\delta q_a(t_A)}\Biggr\}
  \\
  \;&=ie^{iI}\int_{t}^{t'}d\tau \Biggl\{ \text{边界项}-\sum_{b}\frac{\delta q_{b}(\tau) }{\delta q_a(t_A)}\dot{p}_b(\tau)-\frac{\partial H(q(\tau),p(\tau))}{\partial q_b(\tau)}\frac{\delta q_b(\tau)}{\delta q_a(t_A)}\Biggr\}
  \\
  \;&=ie^{iI}\int_{t}^{t'}d\tau \Biggl\{0- \sum_{b}\delta_{ab}\delta(\tau-t_A)\dot{p}_b(\tau)-\frac{\partial H(q(\tau),p(\tau))}{\partial q_b(\tau)}\delta_{ab}\delta(\tau-t_A)\Biggr\}
  \\
  \;&=-ie^{iI} \Biggl\{ \dot{p}_a(t_A)+\frac{\partial H(q(t_A),p(t_A))}{\partial q_a(t_A)}\Biggr\}
\end{aligned}\]
写在一起也就是：
\begin{subequations}\label{eq:can_eq}
  \begin{align}
    & \frac{\delta }{\delta p_a(t_A)}e^{iI}=i\Biggl\{ \dot{q}_a(t_A)-\frac{\partial H(q(t_A),p(t_A))}{\partial p_a(t_A)}\Biggr\}e^{iI}
    \\
    &\frac{\delta }{\delta q_a(t_A)}=-i \Biggl\{ \dot{p}_a(t_A)+\frac{\partial H(q(t_A),p(t_A))}{\partial q_a(t_A)}\Biggr\}e^{iI}
  \end{align}
\end{subequations}
\color{blue}%
在这一组方程中看到了正则方程，上述方程连接了``运动方程''与``全导数''，
说明了``全导数积分为零''有更深一层的物理意义\normalcolor



\medskip\noindent
\color{blue}%
在路径积分中，
$\dot{q}_a(\tau)-\partial H/\partial p_a(\tau)$
在绝大多数路径上\textbf{不为零}
（因为积分变量不受运动方程约束）。
但下面将证明，
对所有路径求和后，
正负贡献精确抵消，期望值为零。
这不是对任意函数都成立的性质
（例如 $q_a^2$ 的期望值一般不为零），
而是运动方程左边所独有的——
因为\textbf{只有它}能被写成 $e^{iI}$ 的全导数。\normalcolor
% =====================================================
%   Schwinger–Dyson 方程（一般形式）
% =====================================================


\medskip\noindent
从路径积分表示 (\ref{eq:PI-field}) 出发，
考虑最一般的编时关联函数，
插入任意多个算符
$\mathcal{O}_A[p(t_A),q(t_A)]$、
$\mathcal{O}_B[p(t_B),q(t_B)]$、$\cdots$。
对 (\ref{eq:IBP-p}) 取被积函数为
$\mathcal{O}_A\,\mathcal{O}_B\,\cdots\,e^{iI}$：
\begin{equation}\label{eq:SD-gen-start}
  0=\int\!\mathcal{D}q\,\frac{\mathcal{D}p}{2\pi}\;
  \frac{\delta}{\delta p_a(t_A)}
  \Bigl[\,
    \mathcal{O}_A\bigl[p(t_A),q(t_A)\bigr]\;
    \mathcal{O}_B\bigl[p(t_B),q(t_B)\bigr]\;\cdots\;
    e^{iI}
  \,\Bigr]\,.
\end{equation}
使用\textbf{泛函莱布尼茨法则}，
$\delta/\delta p_a(t_A)$ 依次作用在每个因子上。
不同时刻的 $p$ 是独立积分变量，
$\delta p_b(t')/\delta p_a(t_A)=\delta_{ab}\,\delta(t'-t_A)$，
$\delta q_b(t')/\delta p_a(t_A)=0$，
因此：
\begin{itemize}
\item 作用在 $\mathcal{O}_A[p(t_A),q(t_A)]$ 上：
  直接给出偏导数 $\partial\mathcal{O}_A/\partial p_a$；
\item 作用在 $\mathcal{O}_B[p(t_B),q(t_B)]$ 上：
  贡献 $(\partial\mathcal{O}_B/\partial p_a)\,\delta(t_A-t_B)$；
\item 作用在后续因子上：类推；
\item 作用在 $e^{iI}$ 上：
  由 (\ref{eq:can_eq}) 给出
  $i(\dot{q}_a-\partial H/\partial p_a)\,e^{iI}$。
\end{itemize}
写成编时期望值并移项，
记全部插入算符为 $\mathcal{O}_J$（$J=A,B,\dots$），
得到\textbf{Schwinger--Dyson 方程}：
\begin{equation}\label{eq:SD-general}
  \biggl\langle T\!\biggl\{
    \biggl(\dot{q}_a(t_A)
      -\frac{\partial H}{\partial p_a(t_A)}\biggr)
    \;\prod_{J}\mathcal{O}_J(t_J)
  \biggr\}\biggr\rangle
  =\;-\,i\;\sum_{K}\delta(t_A-t_K)\;
  \biggl\langle T\!\biggl\{
    \frac{\partial\,\mathcal{O}_K}{\partial\,p_a}\;
    \prod_{J\neq K}\mathcal{O}_J(t_J)
  \biggr\}\biggr\rangle\,.
\end{equation}
对 $q_a(t_A)$ 做泛函导数，类似可得 $p_a$ 运动方程的版本：
\begin{equation}\label{eq:SD-general-q}
  \biggl\langle T\!\biggl\{
    \biggl(\dot{p}_a(t_A)
      +\frac{\partial H}{\partial q_a(t_A)}\biggr)
    \;\prod_{J}\mathcal{O}_J(t_J)
  \biggr\}\biggr\rangle
  =\;-\,i\;\sum_{K}\delta(t_A-t_K)\;
  \biggl\langle T\!\biggl\{
    \frac{\partial\,\mathcal{O}_K}{\partial\,q_a}\;
    \prod_{J\neq K}\mathcal{O}_J(t_J)
  \biggr\}\biggr\rangle\,.
\end{equation}

\medskip\noindent
(\ref{eq:SD-general}) 的右边是对所有插入算符的求和：
第 $K$ 项非零，当且仅当 $\mathcal{O}_K$ 显含 $p_a$，
此时 $\delta(t_A-t_K)$ 将时刻钉死在 $t_K=t_A$——
这称为\textbf{接触项}。
若所有 $\mathcal{O}_J$ 都不含 $p_a$，右边为零，
经典运动方程在编时期望值中精确成立。

\medskip\noindent
在算符形式中，
对编时乘积 $T\{Q_a(t_A)\prod_J\mathcal{O}_J(t_J)\}$
求 $t_A$ 导数时，
$\theta$-函数的导数产生等时对易子：
\begin{equation}\label{eq:T-deriv-gen}
  \frac{\partial}{\partial t_A}\,
  T\Bigl\{Q_a(t_A)\,\prod_{J}\mathcal{O}_J(t_J)\Bigr\}
  =T\Bigl\{\dot{Q}_a(t_A)\,\prod_{J}\mathcal{O}_J\Bigr\}
  +\sum_{K}\delta(t_A-t_K)\;
  T\Bigl\{
    \bigl[Q_a,\,\mathcal{O}_K\bigr]
    \prod_{J\neq K}\mathcal{O}_J
  \Bigr\}\,.
\end{equation}
利用 $[Q_a,\mathcal{O}_K]=i\,\partial\mathcal{O}_K/\partial P_a$
与 Heisenberg 方程 $\dot{Q}_a=\partial H/\partial P_a$，
恰好复现 (\ref{eq:SD-general})。因此
\begin{equation}\label{eq:contact-comm}
  \underbrace{%
    -\,i\;\sum_{K}\delta(t_A-t_K)\;
    \biggl\langle T\!\biggl\{
      \frac{\partial\,\mathcal{O}_K}{\partial\,p_a}\;
      \prod_{J\neq K}\mathcal{O}_J
    \biggr\}\biggr\rangle
  }_{\text{路径积分：接触项}}
  \quad\longleftrightarrow\quad
  \underbrace{%
    \sum_{K}\delta(t_A-t_K)\;
    \bigl\langle T\bigl\{
      [Q_a,\,\mathcal{O}_K]\,
      \textstyle\prod_{J\neq K}\mathcal{O}_J
    \bigr\}\bigr\rangle
  }_{\text{算符形式：等时对易子}}\,.
\end{equation}
路径积分中的接触项与算符形式中的等时对易关系
是同一量子效应的两种语言。


% =============================================================
%   讨论 2：约束系统与黑塞矩阵
% =============================================================
\topictitle{约束系统与黑塞矩阵}

\noindent
\textbf{1.\;从相空间到构型空间}

\medskip\noindent
对许多理论（如实标量场
$\mathcal{L}=\frac{1}{2}\partial_\mu\phi\,\partial^\mu\phi-V(\phi)$），
$p_a$ 在指数 $p\dot{q}-\mathscr{H}$
中仅出现为高斯形式。
对 $p$ 完成高斯积分后，
相空间路径积分退化为构型空间路径积分：
\begin{equation}\label{eq:config-PI}
  Z=\int\!\mathcal{D}q\;e^{iS[q]}\,,\qquad
  S[q]=\int\!d^4x\;\mathcal{L}(q,\partial_\mu q)\,.
\end{equation}
这一步能否执行，取决于\textbf{黑塞矩阵}
\begin{equation}\label{eq:Hessian}
  \mathcal{M}_{ab}
  \equiv
  \frac{\partial^2\mathcal{L}}
    {\partial\dot{q}_a\,\partial\dot{q}_b}
\end{equation}
是否非退化（$\det\mathcal{M}\neq 0$）。
\begin{itemize}
  \item $\det\mathcal{M}\neq 0$：
    Legendre 变换 $p_a=\partial\mathcal{L}/\partial\dot{q}_a$
    可唯一解出 $\dot{q}_a(q,p)$，
    高斯积分产生因子 $(\det\mathcal{M})^{1/2}$
    吸收进测度。
  \item $\det\mathcal{M}=0$：
    不能解出所有速度，
    存在\textbf{约束} $\chi_\alpha(q,p)=0$，
    Legendre 变换失败。
\end{itemize}

\medskip\noindent
\textbf{2.\;约束的物理来源}

\medskip\noindent
以电磁场
$\mathcal{L}=-\frac{1}{4}F_{\mu\nu}F^{\mu\nu}$ 为例，
$\mathcal{L}$ 不含 $\dot{A}_0$，故
$P^0=\partial\mathcal{L}/\partial\dot{A}_0=0$（初级约束）。
黑塞矩阵在 $A_0$ 方向有零本征值——
正是\textbf{规范冗余度}的信号。
对规范理论，Faddeev--Popov 方法在路径积分中
引入规范固定条件和鬼场，
将冗余积分约化为物理自由度上的积分。
这将在非阿贝尔规范理论中系统展开。



% =============================================================
%   讨论 4：欧几里得路径积分
% =============================================================
\topictitle{欧几里得路径积分与非微扰物理}

\noindent
\textbf{1.\;闵可夫斯基路径积分的困难}

\medskip\noindent
(\ref{eq:PI-phase}) 的被积函数
$e^{iI}$（$I\in\mathbb{R}$）的模为
$|e^{iI}|=1$——纯振荡。
这使得数值计算（如格点 QCD）遭遇
符号问题，Monte Carlo 采样失效；
严格分析中收敛性难以定义；
对非微扰效应缺乏数学控制。

\medskip\noindent
\textbf{2.\;Wick 转动}

\medskip\noindent
做解析延拓 $t=-i\tau_E$
（$\tau_E\in\mathbb{R}$，欧几里得时间）。
在构型空间路径积分（已对 $p$ 完成高斯积分后）中，
以 $L=\frac{1}{2}m\dot{q}^2-V(q)$ 为例：
\[
  iS = i\!\int\!dt\,
    \bigl(\tfrac{1}{2}m\dot{q}^2-V\bigr)
  = -\!\int\!d\tau_E\,
    \bigl(\tfrac{1}{2}m\dot{q}_E^{\,2}+V\bigr)
  \equiv -S_E\,,
\]
因此
\begin{equation}\label{eq:Wick}
  e^{\,iS}\;\longrightarrow\;e^{-S_E}\,,\qquad
  S_E\geq 0\,.
\end{equation}
被积函数从纯振荡变为\textbf{指数衰减}：
远离经典路径的构型被强烈压低。
对一般 $3{+}1$ 维场论
（$x^0=-ix^0_E$，度规变为欧几里得），
以实标量场为例：
\begin{equation}\label{eq:ZE}
  Z_E=\int\!\mathcal{D}\phi\;e^{-S_E[\phi]}\,,\qquad
  S_E[\phi]=\int\!d^4x_E\,
    \Bigl(\tfrac{1}{2}(\nabla_E\phi)^2+V(\phi)\Bigr)
    \geq 0\,.
\end{equation}
此形式的优点：
(i)~$e^{-S_E}\leq 1$ 提供自然截断，
收敛性可严格讨论；
(ii)~$e^{-S_E}$ 正定，
Monte Carlo 采样可行——格点场论的基础；
(iii)~$Z_E$ 与统计力学配分函数形式相同
（$S_E\leftrightarrow\beta E$），
重正化群等工具直接移植；
(iv)~解析延拓 $\tau_E\to it$
恢复闵可夫斯基 Feynman 振幅
（Osterwalder--Schrader 重建定理）。

\medskip\noindent
\textbf{3.\;非微扰物理}

\medskip\noindent
路径积分远不止是导出 Feynman 规则的工具。
用路径积分计算欧几里得空间中的振幅，
指数的变量是一个\textbf{负的实量}——
取代了锯齿状路径产生的激烈振荡，
所有远离经典路径的构型被指数式地抑制。
这使得以下非微扰效应可以被系统地处理：

\begin{itemize}
  \item \textbf{瞬子}
  \item \textbf{$\theta$-真空}
  \item \textbf{格点场论与禁闭}
\end{itemize}

\noindent
正如 Weinberg 所强调：
路径积分方法允许我们解释那些
在零耦合常数处有显然奇异性的 S-矩阵贡献，
这样的贡献无法在任何有限阶微扰论中被发现。
自此以后，路径积分方法就变成了
所有使用量子场论的物理学家的工具中
不可分割的一部分。

\subsection{波色子的路径积分}

\subsection{费米子的路径积分}

